\section{The nonunital absorption step}\label{n:sec:absorption}
By \cite[Theorem 4.6 and Definition 4.2]{CETW}, a separable nuclear algebra
with no finite-dimensional quotients, compact nonempty $T$, and uniform
property $\Gamma$ admits unital homomorphisms
\begin{equation}\label{n:eq:mcduff}
 M_n\longrightarrow D\cap A'\qquad(n\ge1).
\end{equation}
The theorem permits nonunital algebras. We verify the passage from
\eqref{n:eq:mcduff} to the nonunital absorption criterion explicitly.

Let
\[
 \Aom=\ell^\infty(A)/\{(a_n):\lim_\omega\|a_n\|=0\},
 \quad J=\ker(\Aom\to D),
 \quad \Fw=(\Aom\cap A')/(\Aom\cap A^\perp).
\]
Here the annihilator consists of the elements whose products on either
side with $A$ are zero.

\begin{lemma}\label{n:lem:central-lift}
The quotient $\Fw$ is unital, the annihilator is contained in $J$, and
the map $\Aom\cap A'\to D\cap A'$ is onto. Consequently, for each $n$
there is an order-zero map $\psi:M_n\to\Aom\cap A'$ satisfying
$\psi(1)+J=1_D$. Denote its composition with the quotient map to $\Fw$
by $\varphi$.
\end{lemma}
\begin{proof}
The approximate-unit sequence defines $u\in\Aom\cap A'$ with
$ua=a=au$ for constant $a\in A$. For $y\in\Aom\cap A'$,
$(uy-y)a=(ua-a)y=0$, and likewise on the left. Thus $u$ is the unit
modulo the annihilator.

If $y$ annihilates $A$, choose a bounded representative $(y_n)$ of norm
at most $M$. For each fixed $k$, $\|y_nu_k\|\to_\omega0$, while
\[
 \|y_n(1-u_k)\|_{2,T}^2
 \le M^2\sup_{\sigma\in T}(1-\sigma(u_k))\longrightarrow0
 \quad(k\to\infty).
\]
Taking these limits in order proves $y\in J$.

For surjectivity, let $v\in\Aom$ lift a member of $D\cap A'$. Its
commutators with $A$ lie in $J$. A quasicentral approximate unit for this
ideal gives positive contractions $d^{(k)}\in J$ that commute with the
first $k$ elements of a dense sequence in $A$ to within $1/k$, and
satisfy $d^{(k)}[v,a]=[v,a]$ to that accuracy on those elements. Choose
positive contraction representatives. On an $\omega$-large set their
uniform $2$-norms are below $1/k$, and their coordinate norm errors are
below $2/k$. The decreasing-set selection used in
Proposition~\ref{n:prop:split} now gives $d\in J\cap A'$ satisfying
$d[v,a]=[v,a]$ for every $a\in A$. Hence $v-dv$ is a central lift of the
prescribed element. This is also the central-surjectivity statement of
\cite[Lemma 1.2]{CETW}.

Apply matrix order-zero lifting to the resulting quotient map and
\eqref{n:eq:mcduff} to obtain $\psi$, and compose with the quotient to
$\Fw$ to obtain $\varphi$. The image of $\psi(1)$ in $D$ is $1_D$.
\end{proof}

A generalized limit trace on $\Aom$ comes from lower-semicontinuous traces
$\nu_n$ on $A$ by
\begin{equation}\label{n:eq:generalized}
 \nu([(b_n)])=\sup_{\epsilon>0}\lim_\omega\nu_n((b_n-\epsilon)_+)
 \quad(b_n\ge0).
\end{equation}
The coordinate traces here need not initially be densely finite. Write
$\Ped(A)$ for the Pedersen ideal, the smallest dense two-sided ideal of
$A$. For
nonzero $a\in\Ped(A)_+$ with $\nu(a)<\infty$, multiplication by $a$
on the norm relative commutant induces a trace $\nu_a$ on $\Fw$.
The relevant largeness condition for $\varphi$ is
$\nu_a(\varphi(1))=\nu(a)$ \cite[Definition 2.4]{CLS}.

We give the normalization argument explicitly under our trace hypotheses;
compare \cite[Lemma 2.3 and the discussion after Definition 2.4]{CLS}.
For a generalized trace finite at a fixed nonzero element of $A_+$,
boundedness expresses the relevant coordinate traces as scalar multiples
of states in $T$, and compactness bounds these scalars uniformly.

\begin{lemma}\label{n:lem:generalized}
Every generalized limit trace $\nu$ finite at a fixed nonzero
$a\in A_+$ is a finite nonnegative multiple of a limit trace through $D$.
The induced maps $\varphi$ in Lemma~\ref{n:lem:central-lift} are therefore
tracially large
for all the required Pedersen-ideal tests.
\end{lemma}
\begin{proof}
Choose $\epsilon>0$ such that $b=(a-\epsilon)_+\ne0$. Every member of
$T$ is faithful, so compactness gives
\[
 m=\min_{\sigma\in T}\sigma(b)>0.
\]
Since $\nu(a)<\infty$, \eqref{n:eq:generalized} gives an $\omega$-large set
on which $\nu_n(b)\le C$ for some finite $C$. A lower-semicontinuous
trace $\eta$ finite on a nonzero positive element of a simple algebra is
necessarily densely finite. Indeed, the set
$N_\eta=\{x:\eta(x^*x)<\infty\}$ is linear by
$(x+y)^*(x+y)\le2x^*x+2y^*y$. Traciality and positivity give
\[
 \eta((xy)^*(xy))=\eta(xyy^*x^*)\le\|y\|^2\eta(x^*x),
 \qquad \eta((yx)^*(yx))\le\|y\|^2\eta(x^*x).
\]
Thus $N_\eta$ is a nonzero two-sided ideal and is norm dense by simplicity.
Approximating square roots shows that finite-trace positives are dense.
Therefore, on this set,
our trace hypothesis gives $\nu_n=r_n\sigma_n$ with $\sigma_n\in T$ and
$0\le r_n\le C/m$. The zero trace is included by $r_n=0$, and the trace
infinite on every nonzero positive element cannot occur there. Extend the
scalars and states boundedly off this set.

For every bounded positive representative $z_n$ and $\delta>0$,
\[
 0\le r_n\sigma_n(z_n)-r_n\sigma_n((z_n-\delta)_+)
       \le(C/m)\delta.
\]
The regularization in \eqref{n:eq:generalized} can therefore be removed,
and
\[
 \nu=(\lim_\omega r_n)(\lim_\omega\sigma_n).
\]
This is a finite, possibly zero multiple of a trace through $D$.
For $z=\psi(1)$ from Lemma~\ref{n:lem:central-lift}, the image of $z$ in
$D$ is $1_D$. Consequently
$\nu_a(\varphi(1))=\nu(az)=\nu(a)$ for every permitted test. The formula
is independent of the representative modulo the annihilator, since that
ideal annihilates $a$.
\end{proof}

The Cuntz semigroup $\Cu(A)$ consists of Cuntz classes of positive elements
of $A\otimes\K$, ordered by subequivalence and added by orthogonal sum.
It is \emph{almost unperforated} if $(k+1)[x]\le k[y]$ implies $[x]\le[y]$
for every integer $k\ge1$.

\begin{proof}[Proof of Theorem~\ref{n:thm:main}]
Theorem~\ref{n:thm:gamma} and \eqref{n:eq:mcduff}, followed by
Lemmas~\ref{n:lem:central-lift} and \ref{n:lem:generalized}, give tracially
large order-zero maps $M_n\to\Fw$ for every $n$. This is uniform
McDuffness in the nonunital sense of \cite[Definition 2.4]{CLS}.

Our comparison hypothesis also makes $\Cu(A)$ almost unperforated.
Indeed, if $(k+1)[x]\le k[y]$ with $y\ne0$, then additivity and
monotonicity of dimension functions give
\[
 d_\tau(x)\le\frac{k}{k+1}<1
 \qquad\text{whenever }d_\tau(y)=1.
\]
The stated comparison gives $[x]\le[y]$. The case $y=0$ is immediate.
This uses precisely the specified finite-target-rank tests, including
the case that this family is empty. Dimension-function monotonicity
follows from cutdown under subequivalence, as in the proof of
Proposition~\ref{n:prop:selection}, followed by increasing the supports
of the cutdowns.

For simple algebras, almost unperforation is the strict comparison
formulation used in \cite[Section 2.3]{CLS}. Its nonunital absorption
criterion \cite[Proposition 5.4]{CLS} states that a simple, separable,
nuclear algebra with this comparison and nontrivial
lower-semicontinuous traces is $\Z$-stable if it is uniformly McDuff.
All these hypotheses have now been verified. Hence $A\cong A\otimes\Z$.
\end{proof}
